\addtocontents{toc}{\protect\Needspace{3\baselineskip}}
\originalchapter{第二次测验复习：轨道稳定性与负压}
公开的 8.286 Fall 2013 第二次测验复习册把 Schwarzschild 圆轨道扰动和活塞功放在同一套题中：前者检验从度规到运动方程的局部稳定性，后者检验从密度变化到总能量的账目。这里的轨道推导也可与本卷引用的既有广义相对论卷对照。

\src{8.286 Quiz 2 Review Problem 18(a)--(d), 原题为2007年 Quiz 2 Q4；官方复习解答。}
\begin{exercise}
令 $h(r)=1-R_S/r$，在 $\theta=\pi/2$ 的 Schwarzschild 度规
$\dd s^2=-hc^2\dd t^2+h^{-1}\dd r^2+r^2\dd\phi^2=-c^2\dd\tau^2$ 中考虑圆轨道 $r=r_0$ 的小径向扰动。
(a) 从 $\frac{\dd}{\dd\tau}(g_{\mu\nu}x'^\nu)=\frac12(\partial_\mu g_{\lambda\sigma})x'^\lambda x'^\sigma$ 求 $r''$，暂可含 $t',r',\phi'$。
(b) 用四速度归一关系消去 $t'^2$，写成 $r''=f_0(r)+f_1(r)\phi'^2$，说明 $r'^2$ 项为何抵消。
(c) 从 $\phi$ 测地线求 $L=r^2\phi'$ 的演化。
(d) 对 $r_0+\delta r$ 线性化，先用 $f_0,f_1$ 给稳定条件，再指出最内稳定圆轨道。
\end{exercise}
\begin{solution}
记撇号为 $\dd/\dd\tau$。径向的欧拉方程是
\[
\frac{\dd}{\dd\tau}(h^{-1}r')
=-\frac12h'c^2t'^2-\frac12h'h^{-2}r'^2+r\phi'^2,
\]
即 $r''=-\frac12hh'c^2t'^2+\frac12(h'/h)r'^2+hr\phi'^2$。
归一式给 $hc^2t'^2=c^2+h^{-1}r'^2+r^2\phi'^2$；代入后两处 $r'^2$ 项恰好抵消，
\[
r''=-\frac12c^2h'+(hr-\tfrac12r^2h')\phi'^2
=-\frac{GM}{r^2}+
\left(r-\frac{3R_S}{2}\right)\phi'^2 ,
\quad R_S=\frac{2GM}{c^2}.
\]
故 $f_0=-GM/r^2,\ f_1=r-3R_S/2$。由于度规不含 $\phi$，$L=r^2\phi'$ 守恒。固定此 $L$ 定义
$F(r)=f_0(r)+f_1(r)L^2/r^4$，圆轨道满足 $F(r_0)=0$，小扰动满足
$\delta r''=F'(r_0)\delta r$；$F'(r_0)<0$ 才发生有界谐振，等号为临界。显式代入可得存在类时圆轨道须 $r_0>3R_S/2$，稳定须 $r_0>3R_S$。$3R_S/2<r_0<3R_S$ 有圆轨道但不稳定，$r_0=3R_S$ 为最内稳定边界；本问测的是固定角动量扰动，不是只凭圆轨道存在性下结论。
\end{solution}

\src{8.286 Quiz 2 Review Problem 19(a)--(c), 原题为2002年 Quiz 3 Q3；官方复习解答。}
\begin{exercise}
一均匀“奇异物质”在绝热活塞中密度 $u\propto V^{-1/2}$，初始体积 $V$、能量密度 $u_0$。活塞拉大无穷小体积 $\Delta V$：(a) 求内部总能量变化 $\Delta U$，增加取正；(b) 压强为 $p$ 时求外界\emph{对系统}所做的功 $\Delta W$；(c) 求 $p$ 与 $u$ 的关系。此指数和正式 Quiz 2 Q4 的 $V^{-3/2}$ 不同。
\end{exercise}
\begin{solution}
写 $u(V')=u_0(V/V')^{1/2}$，因此 $U(V')=u(V')V'=u_0V^{1/2}(V')^{1/2}$，微分得
$\Delta U=\tfrac12u_0\Delta V$。外界对系统所做的压缩功按符号约定是 $\Delta W=-p\Delta V$。绝热且没有其他功、热项时 $\Delta U=\Delta W$，于是 $p=-u/2$。负压使拉大体积时系统\emph{总}能量增加，和密度下降并不矛盾；密度与总体积同时变化。一般若 $u\propto V^{-(1+w)}$，上述同一预算给 $p=wu$。
\end{solution}
